\section{Discussion and Limitations}
\label{sec:discussion}

\paragraph{What produced the accuracy.}
The largest single jump came from the training recipe rather than the
architecture: the 95.42\% flip-TTA baseline became 97.36\% with long training,
TrivialAugment, Cutout, mixup/CutMix, SAM, and EMA. KD at 600 epochs reaches
97.74\%, but that comparison also changes the schedule, so it is an upper bound
on the KD effect. The wider student adds 0.24 points at 300 epochs but only 0.08
on the 600-epoch main path, within noise. The 4.88M student ends up better than
the 28.5M PyramidNet teacher (97.64\%) and within 0.3 points of the two-teacher
pool (98.14\%), but it does not beat either 36.5M WRN-28-10 teacher
(98.06/98.12). Architecture search inside the 5M budget was secondary.

\paragraph{The unexplained gap.}
A4 reaches 97.92\% with 5-crop and V2 reaches 97.94\%, a 0.02 difference that is
not meaningful; we treat the 97.9\% range as the headline holdout result. The
selected submissions scored 0.911 and 0.904 public, and the team's private score
was 0.880, 13th of 81. The private score is roughly ten points below the holdout
range, and the gap between the best public score (0.911) and the private score
(0.880) is 3.1 points. Because the public and private halves are 1,000 images
each, part of that difference is sampling noise; the rest could be a harder
private subset, a different capture distribution, or both. We do not have the
final evaluation labels, and we do not analyze the provided unlabeled images
here, so we cannot separate these causes.

\paragraph{Why one-view AugMix was not shipped.}
One-view AugMix is the one robustness intervention that survived every check we
ran: it is clean-neutral at 300 epochs, improves the composite by 2.2 points,
improves photographic jitter by 1.1--1.9 points, replicates across seeds, and
keeps its effect under three-teacher KD. It was still not in the final submission.
The KD confirmation had not finished when the team selected its submissions; its
logs stop at epoch 178/300, and the full-data 1,200-epoch AugMix run was never
trained. Given that the final private gap is unexplained, we cannot claim that
shipping AugMix would have improved the score. The honest conclusion is that it
was the best-supported robustness candidate we had, and the submission schedule
did not allow us to test it at the shipping scale.

\paragraph{Limitations.}
Most comparisons use a single seed; the exceptions are the two AugMix seeds and
the ASAM confirm. The holdout is 5,000 images, so the seed-to-seed spread is
0.2--0.4 points and small differences should not be over-interpreted. The B5 KD
comparison is provisional at 200/300 epochs and was interrupted twice; its
partial curves are in Figure~\ref{fig:curves}(f). The proxies are synthetic and label-free stress tests, not models of the final
shift. Both leaderboards are 1,000 images, so they are noisy. We have no completed fresh 1,200-epoch ASAM,
SWA, or LS-KD run (F1200\_asam stopped at epoch 250), and we did not test DKD or
structural teacher pruning. Finally, the full-data models have no unbiased
validation score, so the final selection relies on the holdout models and, for
the full-data pair, the public leaderboard.

\paragraph{What we would do differently.}
First, ablate augmentation at the shipping schedule: the 150-epoch pilot
overstated AugMix's effect because the 300-epoch control improved so much. Second, keep per-epoch logs for every run in the repository; several
curves in this paper are missing because the logs stayed on a cluster or a cloud
volume. Third, complete the KD robustness run before submission selection, so
that a clean-neutral augmentation has a chance to ship. Fourth, run the main
comparisons with at least two seeds; the single-seed spread is large enough to
swallow several of the effects we measured.
